\documentclass[12pt]{article}
\usepackage[a4paper,margin=12mm]{geometry}
\usepackage{tikz}
\usetikzlibrary{calc}
\pagestyle{empty}

\begin{document}
\begin{tikzpicture}[remember picture,overlay]
% ======= Adjustable layout parameters (all in mm) =======
\def\Margin{12mm}

\def\CornerSize{5mm}
\def\CornerOffset{5mm}

\def\HeaderH{75mm}

\def\BubbleD{4mm}
\def\RowPitch{6mm}
\def\ColPitch{8mm}

\def\GrpPitch{12mm}
\def\VerPitch{14mm}

\def\LabelPad{4mm}
\def\BoxGap{6mm}
\def\BoxStroke{0.6mm}
\def\Stroke{0.4mm}

% ======= Page anchors =======
\coordinate (NW) at (current page.north west);
\coordinate (NE) at (current page.north east);
\coordinate (SW) at (current page.south west);
\coordinate (SE) at (current page.south east);

% Inner content frame (margins)
\coordinate (TL) at ($(NW)+(\Margin,-\Margin)$);
\coordinate (TR) at ($(NE)+(-\Margin,-\Margin)$);
\coordinate (BL) at ($(SW)+(\Margin,\Margin)$);
\coordinate (BR) at ($(SE)+(-\Margin,\Margin)$);

% ======= Corner marks (solid squares) =======
\draw[fill=black] ($(NW)+(\CornerOffset,-\CornerOffset)$) rectangle ++(\CornerSize,-\CornerSize);
\draw[fill=black] ($(NE)+(-\CornerOffset-\CornerSize,-\CornerOffset)$) rectangle ++(\CornerSize,-\CornerSize);
\draw[fill=black] ($(SW)+(\CornerOffset,\CornerOffset+\CornerSize)$) rectangle ++(\CornerSize,-\CornerSize);
\draw[fill=black] ($(SE)+(-\CornerOffset-\CornerSize,\CornerOffset+\CornerSize)$) rectangle ++(\CornerSize,-\CornerSize);

% ======= Title & meta (top-right inside header, 3 lines) =======
\node[anchor=north east] at ($(TR)+(0,2mm)$) {%
  \shortstack[r]{\Large
    Mathematics for Business and Economics I \\[3pt]
    Prof. Vincent Ginis, Melika Mobini \& Vincent Holst \\[3pt]
    Bonus Test 
  }
};

% ======= Student ID OMR (7x10) =======
\node[anchor=north west] at ($(TL)+(18mm,2mm)$) {\bfseries Student ID};

% Grid origin (center of top-left bubble), relative to TL
\def\IDLeft{8mm}  % x-offset from TL
\def\IDTop{8mm}  % y-offset from TL downward

\def\IDCols{7}
\def\IDRows{10}

% Row labels 0..9 (left of the grid)
\foreach \r in {0,...,9}{
  \node[anchor=east,inner sep=0pt] at ($(TL)+(\IDLeft-5mm,-\IDTop - \r*\RowPitch)$) {\footnotesize \r};
}

% Bubbles: 7 columns x 10 rows
\foreach \r in {0,...,9}{
  \foreach \c in {0,...,6}{
    \draw[line width=\Stroke] ($(TL)+(\IDLeft+\c*\ColPitch,-\IDTop-\r*\RowPitch)$) circle (\BubbleD/2);
  }
}

% ======= Group OMR (1–5), to the right of ID grid =======
\def\GrpLeft{83mm}  
\def\GrpTop{28mm}   

\node[anchor=south] at ($(TL)+(\GrpLeft+2*\GrpPitch,-\GrpTop+\BubbleD/2+3mm)$) {\bfseries Group};
\foreach \i in {0,...,4}{
  \draw[line width=\Stroke] ($(TL)+(\GrpLeft+\i*\GrpPitch,-\GrpTop)$) circle (\BubbleD/2);
  \node[anchor=north,inner sep=0pt] at ($(TL)+(\GrpLeft+\i*\GrpPitch,-\GrpTop-4.5mm)$) {\footnotesize \the\numexpr\i+1\relax};
}

% ======= Version OMR (A/B) =======
\def\VerLeft{\GrpLeft}
\def\VerTop{54mm} 

% Label centered above the version bubbles
\node[anchor=south] at ($(TL)+(\VerLeft+0.5*\VerPitch,-\VerTop+\BubbleD/2+3mm)$) {\bfseries Version};

% Bubbles + labels
\foreach \i/\lab in {0/A,1/B}{
  \coordinate (VerBubble) at ($(TL)+(\VerLeft+\i*\VerPitch,-\VerTop)$);
  \draw[line width=\Stroke] (VerBubble) circle (\BubbleD/2);
  \node[anchor=north,inner sep=0pt] at ($(VerBubble)+(0,-4.5mm)$) {\footnotesize \lab};
}

\node[anchor=north west, text=black, align=left, text width=0.42\textwidth] 
  at ($(TR)+(-80mm,-38 mm)$) {%
    Disclaimer: 
    \footnotesize Only answers in the black boxes will be graded. 
    Write down intermediate steps for full points. 
    If your student ID ends with an even number, 
    make version A. Otherwise, make version B. 
    Indicate your group and version left.
};

% ======= Answer regions =======
\def\Q1H{70mm}
\def\BodyTop{\HeaderH+6mm}

% Q1 box
\draw[line width=\BoxStroke] ($(TL)+(0,-\BodyTop)$) rectangle ($(TR)+(0,-\BodyTop-\Q1H)$);
\node[anchor=north west,inner sep=2mm] at ($(TL)+(0,-\BodyTop)$) {\bfseries Question 1};

% ROI ticks for Q1
\def\Tick{4mm}
\coordinate (Q1TL) at ($(TL)+(0,-\BodyTop)$);
\coordinate (Q1TR) at ($(TR)+(0,-\BodyTop)$);
\coordinate (Q1BL) at ($(TL)+(0,-\BodyTop-\Q1H)$);
\coordinate (Q1BR) at ($(TR)+(0,-\BodyTop-\Q1H)$);
\draw[line width=\Stroke] (Q1TL) -- ++(\Tick,0) (Q1TL) -- ++(0,-\Tick);
\draw[line width=\Stroke] (Q1TR) -- ++(-\Tick,0) (Q1TR) -- ++(0,-\Tick);
\draw[line width=\Stroke] (Q1BL) -- ++(\Tick,0) (Q1BL) -- ++(0,\Tick);
\draw[line width=\Stroke] (Q1BR) -- ++(-\Tick,0) (Q1BR) -- ++(0,\Tick);

% Q2 box
\def\QGap{6mm}
\draw[line width=\BoxStroke] ($(TL)+(0,-\BodyTop-\Q1H-\QGap)$) rectangle ($(BR)$);
\node[anchor=north west,inner sep=2mm] at ($(TL)+(0,-\BodyTop-\Q1H-\QGap)$) {\bfseries Question 2};

% ROI ticks for Q2
\coordinate (Q2TL) at ($(TL)+(0,-\BodyTop-\Q1H-\QGap)$);
\coordinate (Q2TR) at ($(TR)+(0,-\BodyTop-\Q1H-\QGap)$);
\coordinate (Q2BL) at ($(BL)$);
\coordinate (Q2BR) at ($(BR)$);
\draw[line width=\Stroke] (Q2TL) -- ++(\Tick,0) (Q2TL) -- ++(0,-\Tick);
\draw[line width=\Stroke] (Q2TR) -- ++(-\Tick,0) (Q2TR) -- ++(0,-\Tick);
\draw[line width=\Stroke] (Q2BL) -- ++(\Tick,0) (Q2BL) -- ++(0,\Tick);
\draw[line width=\Stroke] (Q2BR) -- ++(-\Tick,0) (Q2BR) -- ++(0,\Tick);
\end{tikzpicture}
\end{document}
